% ============================================================================
% METHODOLOGY: Key Formulas in LaTeX
% FTLE Calculation, Perturbation Site Selection, Cloud Seeding Simulation
% ============================================================================

\documentclass{article}
\usepackage{amsmath}
\usepackage{amssymb}
\usepackage{bm}

\begin{document}

% ============================================================================
\section*{1. FTLE (Finite-Time Lyapunov Exponent)}
% ============================================================================

\subsection*{1.1 Particle Advection (Euler Integration)}

At each time step $n$, advect particles using wind field:

\begin{align}
\text{lat}_{n+1} &= \text{lat}_n + \frac{v \cdot \Delta t}{111{,}000 \text{ m}} \\
\text{lon}_{n+1} &= \text{lon}_n + \frac{u \cdot \Delta t}{111{,}000 \text{ m} \cdot \cos(\text{lat}_n)}
\end{align}

where:
\begin{itemize}
    \item $u, v$ = zonal and meridional wind components (m/s)
    \item $\Delta t$ = time step (21,600 s = 6 hours)
    \item $111{,}000$ m $\approx$ 111 km/degree
\end{itemize}

% ----------------------------------------------------------------------------
\subsection*{1.2 Flow Map Gradient (Jacobian Matrix)}

Deformation gradient tensor computed using centered finite differences:

\begin{equation}
\frac{\partial \bm{\Phi}}{\partial \bm{x}_0} =
\begin{bmatrix}
\dfrac{\partial \text{lat}_f}{\partial \text{lat}_0} & \dfrac{\partial \text{lat}_f}{\partial \text{lon}_0} \\[0.5em]
\dfrac{\partial \text{lon}_f}{\partial \text{lat}_0} & \dfrac{\partial \text{lon}_f}{\partial \text{lon}_0}
\end{bmatrix}
\end{equation}

Finite difference approximations:
\begin{align}
\frac{\partial \text{lat}_f}{\partial \text{lat}_0} &\approx \frac{\text{lat}_f[i+1,j] - \text{lat}_f[i-1,j]}{2 \cdot \Delta \text{lat}} \\
\frac{\partial \text{lat}_f}{\partial \text{lon}_0} &\approx \frac{\text{lat}_f[i,j+1] - \text{lat}_f[i,j-1]}{2 \cdot \Delta \text{lon} \cdot \cos(\text{lat})} \\
\frac{\partial \text{lon}_f}{\partial \text{lat}_0} &\approx \frac{\text{lon}_f[i+1,j] - \text{lon}_f[i-1,j]}{2 \cdot \Delta \text{lat}} \\
\frac{\partial \text{lon}_f}{\partial \text{lon}_0} &\approx \frac{\text{lon}_f[i,j+1] - \text{lon}_f[i,j-1]}{2 \cdot \Delta \text{lon} \cdot \cos(\text{lat})}
\end{align}

% ----------------------------------------------------------------------------
\subsection*{1.3 Cauchy-Green Strain Tensor}

\begin{equation}
\bm{C} = \left(\frac{\partial \bm{\Phi}}{\partial \bm{x}_0}\right)^T \cdot \left(\frac{\partial \bm{\Phi}}{\partial \bm{x}_0}\right)
\end{equation}

% ----------------------------------------------------------------------------
\subsection*{1.4 FTLE Calculation}

\begin{equation}
\boxed{
\text{FTLE} = \frac{1}{2T} \ln(\lambda_{\max})
}
\end{equation}

where:
\begin{itemize}
    \item $\lambda_{\max}$ = largest eigenvalue of Cauchy-Green tensor $\bm{C}$
    \item $T$ = integration time (48 hours = 2 days for tropical cyclones)
    \item Units: day$^{-1}$
\end{itemize}

Alternative formulation (using singular values):
\begin{equation}
\text{FTLE} = \frac{1}{T} \ln\left(\sqrt{\lambda_{\max}}\right) = \frac{1}{T} \ln(\sigma_{\max})
\end{equation}

where $\sigma_{\max}$ is the maximum singular value of $\frac{\partial \bm{\Phi}}{\partial \bm{x}_0}$.


% ============================================================================
\section*{2. Perturbation Site Selection}
% ============================================================================

\subsection*{2.1 FTLE Threshold (Sensitivity Criterion)}

Select top $p$-th percentile of FTLE values:

\begin{equation}
\text{FTLE}_{\text{threshold}} = P_p(\text{FTLE field})
\end{equation}

where $P_p$ denotes the $p$-th percentile function (typically $p = 85$, i.e., top 15\%).

Candidate set:
\begin{equation}
\mathcal{S}_{\text{FTLE}} = \left\{ (\text{lat}, \text{lon}) \,|\, \text{FTLE}(\text{lat}, \text{lon}) \geq \text{FTLE}_{\text{threshold}} \right\}
\end{equation}

% ----------------------------------------------------------------------------
\subsection*{2.2 Distance from TC Center (Environmental Steering Criterion)}

Haversine distance formula:
\begin{align}
a &= \sin^2\left(\frac{\Delta \text{lat}}{2}\right) + \cos(\text{lat}_1) \cdot \cos(\text{lat}_2) \cdot \sin^2\left(\frac{\Delta \text{lon}}{2}\right) \\
c &= 2 \cdot \arcsin(\sqrt{a}) \\
d &= R_{\text{Earth}} \cdot c
\end{align}

where:
\begin{itemize}
    \item $\Delta \text{lat} = \text{lat}_2 - \text{lat}_1$, $\Delta \text{lon} = \text{lon}_2 - \text{lon}_1$
    \item $R_{\text{Earth}} = 6371$ km
\end{itemize}

Environmental steering region filter:
\begin{equation}
\mathcal{S}_{\text{dist}} = \left\{ (\text{lat}, \text{lon}) \in \mathcal{S}_{\text{FTLE}} \,|\, d_{\min} \leq d(\text{lat}, \text{lon}, \text{lat}_{\text{TC}}, \text{lon}_{\text{TC}}) \leq d_{\max} \right\}
\end{equation}

where typically $d_{\min} = 500$ km and $d_{\max} = 1500$ km for TC track modification.

% ----------------------------------------------------------------------------
\subsection*{2.3 Minimum Site Separation}

For any two selected sites $i$ and $j$:
\begin{equation}
d(\text{lat}_i, \text{lon}_i, \text{lat}_j, \text{lon}_j) \geq d_{\text{sep}}
\end{equation}

where $d_{\text{sep}} = 300$ km (avoids redundant nearby sites).

% ----------------------------------------------------------------------------
\subsection*{2.4 Seeding Mask Creation}

For each selected site center $(\text{lat}_c, \text{lon}_c)$ with radius $R = 300$ km:

\begin{equation}
M[i,j] =
\begin{cases}
1 & \text{if } d(\text{lat}[i], \text{lon}[j], \text{lat}_c, \text{lon}_c) \leq R \\
0 & \text{otherwise}
\end{cases}
\end{equation}

Combined mask for $N$ sites:
\begin{equation}
\boxed{
M_{\text{total}} = \bigvee_{k=1}^{N} M_k = \text{OR}(M_1, M_2, \ldots, M_N)
}
\end{equation}


% ============================================================================
\section*{3. Cloud Seeding Perturbation Simulation}
% ============================================================================

\subsection*{3.1 Saturation Mixing Ratio (Clausius-Clapeyron)}

Saturation vapor pressure:
\begin{equation}
e_{\text{sat}}(T) = 611.2 \text{ Pa} \cdot \exp\left[\frac{L_v}{R_v} \left(\frac{1}{T_0} - \frac{1}{T}\right)\right]
\end{equation}

Saturation mixing ratio:
\begin{equation}
q_{\text{sat}} = \frac{\varepsilon \cdot e_{\text{sat}}}{P - e_{\text{sat}}}
\end{equation}

where:
\begin{itemize}
    \item $L_v = 2.5 \times 10^6$ J/kg (latent heat of vaporization)
    \item $R_v = 461.5$ J/(kg·K) (gas constant for water vapor)
    \item $\varepsilon = 0.622$ (molecular weight ratio $M_d / M_v$)
    \item $T_0 = 273.15$ K
    \item $P$ = pressure (Pa)
\end{itemize}

% ----------------------------------------------------------------------------
\subsection*{3.2 Relative Humidity}

\begin{equation}
\text{RH} = \frac{q}{q_{\text{sat}}}
\end{equation}

where $q$ is the specific humidity (kg/kg).

% ----------------------------------------------------------------------------
\subsection*{3.3 Water Vapor Frozen (Ice Nucleation/Deposition)}

Amount of vapor frozen via deposition process:
\begin{equation}
q_{\text{frozen}} = q_{\text{old}} \cdot \eta_{\text{freeze}} \cdot \text{RH} \cdot M
\end{equation}

where:
\begin{itemize}
    \item $\eta_{\text{freeze}} = 0.60$ (freeze efficiency, tunable parameter)
    \item $M$ = spatial seeding mask
\end{itemize}

% ----------------------------------------------------------------------------
\subsection*{3.4 Precipitation Removal}

Precipitation fallout:
\begin{equation}
q_{\text{precip}} = q_{\text{frozen}} \cdot f_{\text{fallout}}
\end{equation}

Actual moisture removed (with safety cap):
\begin{equation}
q_{\text{removed}} = \min(q_{\text{precip}}, \, q_{\text{old}} \cdot f_{\max})
\end{equation}

where:
\begin{itemize}
    \item $f_{\text{fallout}} = 0.80$ (80\% precipitates out)
    \item $f_{\max} = 0.50$ (maximum 50\% of available moisture, safety limit)
\end{itemize}

% ----------------------------------------------------------------------------
\subsection*{3.5 Latent Heat Release}

Energy released during deposition (vapor $\rightarrow$ ice):
\begin{equation}
Q_{\text{released}} = q_{\text{removed}} \cdot L_d
\end{equation}

Temperature change:
\begin{equation}
\Delta T = \frac{Q_{\text{released}}}{C_p}
\end{equation}

where:
\begin{itemize}
    \item $L_d = L_v + L_f = 2.834 \times 10^6$ J/kg (deposition latent heat)
    \item $L_f = 3.34 \times 10^5$ J/kg (latent heat of fusion)
    \item $C_p = 1004$ J/(kg·K) (specific heat of air at constant pressure)
\end{itemize}

% ----------------------------------------------------------------------------
\subsection*{3.6 Apply Perturbations}

New atmospheric state:
\begin{equation}
\boxed{
\begin{aligned}
T_{\text{new}} &= T_{\text{old}} + \Delta T \cdot M \\
q_{\text{new}} &= q_{\text{old}} - q_{\text{removed}} \cdot M
\end{aligned}
}
\end{equation}

where $M$ is the seeding mask applied at pressure levels $P = [700, 500, 300]$ hPa.

% ----------------------------------------------------------------------------
\subsection*{3.7 Physical Validation}

After applying perturbations, recalculate saturation and validate bounds:

\begin{align}
q_{\text{sat}}^{\text{new}} &= f(T_{\text{new}}, P) \quad \text{(recalculate with new temperature)} \\
\text{RH}^{\text{new}} &= \frac{q_{\text{new}}}{q_{\text{sat}}^{\text{new}}} \\
0 &\leq \text{RH}^{\text{new}} \leq 1 \quad \text{(enforce bounds)} \\
q_{\text{new}} &\geq 0 \quad \text{(no negative moisture)}
\end{align}


% ============================================================================
\section*{Summary of Key Parameters (Sandy 2012 Case)}
% ============================================================================

\begin{table}[h]
\centering
\begin{tabular}{lll}
\hline
\textbf{Parameter} & \textbf{Symbol} & \textbf{Value} \\
\hline
FTLE integration time & $T$ & 48 hours \\
FTLE percentile threshold & $p$ & 85\% \\
TC distance range & $[d_{\min}, d_{\max}]$ & [500, 1500] km \\
Site radius & $R$ & 300 km \\
Site separation & $d_{\text{sep}}$ & 300 km \\
Pressure levels & $P$ & [700, 500, 300] hPa \\
Freeze efficiency & $\eta_{\text{freeze}}$ & 0.60 \\
Fallout fraction & $f_{\text{fallout}}$ & 0.80 \\
Max removal fraction & $f_{\max}$ & 0.50 \\
Deposition latent heat & $L_d$ & $2.834 \times 10^6$ J/kg \\
Specific heat of air & $C_p$ & 1004 J/(kg·K) \\
\hline
\end{tabular}
\end{table}

\subsection*{Typical Perturbation Magnitudes}

\begin{align}
\Delta T &: +0.5 \text{ to } +3.0 \text{ K} \quad \text{(warming from latent heat)} \\
\Delta q &: -0.5 \text{ to } -2.0 \text{ g/kg} \quad \text{(drying from precipitation)} \\
Q_{\text{released}} &: \sim 1\text{--}3 \text{ MJ/kg} \\
\text{Seeded area} &: \sim 850{,}000 \text{ km}^2 \quad \text{(985 grid cells at 0.25° resolution)} \\
\text{Seeded mass} &: \sim 10^{12}\text{--}10^{13} \text{ kg}
\end{align}

\end{document}
